\documentclass[11pt,letterpaper]{article}
\usepackage[spanish,es-noshorthands]{babel}
\usepackage[T1]{fontenc}
\usepackage{lmodern}
\usepackage[margin=2.3cm]{geometry}
\usepackage{amsmath}
\usepackage[hidelinks]{hyperref}
\setlength{\parindent}{0pt}
\setlength{\parskip}{5pt}
\begin{document}
{\LARGE\bfseries Título del tema\par}
Cálculo Diferencial e Integral · Unidad y versión\par
\section{Teoría}
\textbf{Objetivo.} Escriba el aprendizaje esperado.

\textbf{Conocimientos previos.} Identifique los conocimientos que requiere el tema.

\subsection{Conceptos y notación}
Explique las definiciones, su notación y las condiciones de validez disponibles.

\section{Ejemplos resueltos}
\subsection{Ejemplo 1. Título}
\textbf{Enunciado.} Incluya el dominio cuando corresponda.

\textbf{Análisis.} Explique por qué se elige el procedimiento.
\begin{enumerate}
\item \textbf{Acción del primer paso.} Desarrolle el paso.
\end{enumerate}
\textbf{Resultado.} Indique la respuesta y la comprobación disponible.

\section{Práctica}
\subsection{Ejercicio 1. Título}
\label{ej:p1}
Transcriba el enunciado y el dominio. Consulte la solución en la sección~\ref{sol:p1}.

\section{Soluciones}
\subsection{Solución del ejercicio 1}
\label{sol:p1}
\textbf{Análisis.} Explique el procedimiento.
\begin{enumerate}
\item \textbf{Acción del primer paso.} Desarrolle el paso.
\end{enumerate}
\textbf{Resultado.} Indique la respuesta y la comprobación disponible.

Volver al ejercicio~\ref{ej:p1}.
\end{document}
